\section{Method}
\label{sec:method}

\begin{figure*}[t]
\centering
\resizebox{\textwidth}{!}{% Fig. 2 (one frame of VSMT-lean), TikZ; included by sections/method.tex. Drawn by hand from docs/METHOD.md sections
% 3-8; no data. Frozen front end -> recall and sealed features -> joint assignment -> existence decisions -> executor;
% the teacher (training only) labels the sealed rows and columns; bottom right: the entity states (the map is the set
% of active and dormant entities; a retracted entity is out of the map but stays recallable). The cup example of
% Fig. 1 is not repeated here.
\begin{tikzpicture}[
  font=\scriptsize, >={Stealth[length=3.5pt]},
  box/.style={draw, rounded corners=1.5pt, align=center, inner sep=2pt, fill=white, minimum height=12mm},
  learned/.style={draw=orange!70!black, fill=orange!12, rounded corners=1.5pt, align=center, inner sep=2pt, minimum height=6mm},
  private/.style={draw=blue!55!black, fill=blue!6, dashed, rounded corners=1.5pt, align=center, inner sep=2pt, minimum height=6mm},
  seal/.style={fill=black!75, text=white, rounded corners=1pt, inner sep=1.2pt, font=\fontsize{6}{6.6}\selectfont\bfseries},
  lab/.style={font=\fontsize{6}{6.6}\selectfont, align=center, inner sep=0.8pt},
  st/.style={draw, rounded corners=1pt, font=\scriptsize, inner sep=1.1pt, minimum height=3.6mm, fill=white},
]
\node[box, fill=gray!10, text width=28mm] (fe) at (15mm, 0)
  {\textbf{Frozen front end}\\RGB-D frame $\to$ anonymous fragments with descriptors, 3D geometry; free space};
\node[box, text width=26mm] (rec) at (51mm, 0)
  {\textbf{Recall}\\for each fragment: 5 most similar entities within 3\,m $\cup$ 3 most similar anywhere};
\node[box, text width=32mm] (asg) at (91mm, 0)
  {\textbf{Joint assignment}\\fragments $\times$ (recalled entities + one \op{BIRTH} column each), one minimum-cost solve $\to$ \op{BIND} / \op{REACTIVATE} / \op{BIRTH}};
\node[box, text width=27mm] (ex) at (131.5mm, 0)
  {\textbf{Existence decision}\\each unmatched entity that should be visible $\to$ \op{RETRACT} or \op{NOOP}};
\node[box, text width=25mm] (exe) at (166.5mm, 0)
  {\textbf{Executor}\\applies the whole program to a copy of $M_{t-1}$; all or nothing; history kept};
\draw[->] (fe) -- (rec);
\draw[->] (rec) -- node[seal, above=1mm] {seal A} (asg);
\draw[->] (asg) -- node[seal, above=1mm] {seal B} (ex);
\draw[->] (ex) -- (exe);
\draw[->] (exe.north) -- ++(0, 2.5mm) -| node[lab, pos=0.25, above] {$M_t$, the memory of the next frame} (rec.north);
\node[learned, text width=54mm] (heads) at (107mm, -13.6mm)
  {\textbf{The only learned part:} cost heads $a(f,e)$ (association), $b(f)$ (birth), $r(e)$ (existence); 54,787 parameters};
\draw[->, orange!70!black] (heads.north -| asg) -- (asg.south);
\draw[->, orange!70!black] (heads.north -| ex) -- (ex.south);
\node[private, text width=50mm] (teach) at (37mm, -13.6mm)
  {\textbf{Training and evaluation only:} the teacher reads private instance truth and labels only the sealed rows and columns};
\draw[->, dashed, blue!55!black] (teach) -- (heads);
% entity states
\draw[black!45, dashed, rounded corners=1.5pt] (138.5mm, -7.9mm) rectangle (161.5mm, -20.4mm);
\node[font=\fontsize{6}{6.6}\selectfont, text=black!60, anchor=north east, inner sep=0.6pt] at (161.2mm, -8.2mm) {map};
\node[st] (sA) at (150mm, -11.1mm) {active};
\node[st] (sD) at (150mm, -17.3mm) {dormant};
\draw[->] ([xshift=-1.6mm]sA.south) -- node[lab, left=0.4mm] {3 misses} ([xshift=-1.6mm]sD.north);
\draw[->] ([xshift=1.6mm]sD.north) -- node[lab, right=0.4mm] {match} ([xshift=1.6mm]sA.south);
\node[st, fill=black!7, dashed] (sR) at (173.5mm, -14.15mm) {retracted};
\draw[->] (161.5mm, -12.1mm) -- node[lab, above] {\op{RETRACT}} (sR.north west);
\draw[->, dashed] (sR.south west) -- node[lab, below=0.6mm] {\op{REACTIVATE}} (161.5mm, -16.2mm);
\end{tikzpicture}
}
\caption{(a) One frame. Recall sets and association and birth features are sealed before any private file is opened (A), assignments and existence features after the solve (B); the teacher labels only sealed rows and columns, in training only. Orange: the only learned part. (b) A cup moved out of view: VSMT-lean reactivates the same identity, \arm{NoVersion} and \arm{AssocOnly} create a new one, and \arm{AssocOnly} also keeps a stale entity. Illustration, not data.}
\label{fig:method}
\end{figure*}

\subsection{Problem, memory and the five atoms}
Fig.~\ref{fig:method} summarises one frame and the lifecycle of a moved object.
At frame $t$ the memory receives public observations $O_t$ from a frozen front end---anonymous mask fragments with descriptors and 3D geometry, free-space and visibility evidence from depth, and the causal pose---and its own previous memory $M_{t-1}$; it outputs a frame program $\Pi_t$ and $M_t=\mathrm{exec}(M_{t-1},\Pi_t)$.
No simulator object identity, ground-truth mask, intervention log or future observation is an input; the executor is fixed, and what is learned is how $\Pi_t$ is chosen.
The memory is a set of entity records holding an anonymous identity, a state in $\{\text{active},\text{dormant},\text{retracted}\}$, a version chain (predecessor, opening and closing frame, closing transaction), a running-mean and a best-view descriptor, the centroid and box of the surface seen in the most recent frame, observation and missed-opportunity counts, and the evidence list.
\op{NOOP} keeps an unmatched, should-be-visible entity and increments its missed-opportunity count; \op{BIND} attaches a fragment to an active entity; \op{BIRTH} creates an entity from an unassigned fragment; \op{RETRACT} closes the current version of an unmatched, should-be-visible entity whose existence head fires, keeping its evidence and history; \op{REACTIVATE} opens a new version of a dormant or retracted entity at the new place, with the same identity.
A replacement is the composite \op{RETRACT}+\op{BIRTH}.
A shared rule moves an active entity to dormant (``not seen for a while, no evidence of absence'') after three consecutive missed opportunities and back on any match (\arm{AssocOnly} has no dormant state); retracted means ``evidence that it is no longer there''.
A shared deduplication runs every 10 frames over active and dormant entities (cosine $\geq 0.6$, centroid distance $\leq 0.5$\,m, box IoU $\geq 0.05$) and folds one record into the other, keeping all evidence.
The executor applies the whole program on a clone of $M_{t-1}$, checks every precondition and rolls the frame back if any atom fails; history is never physically deleted.

\subsection{Frozen front end}
Fragments come from the simulator's per-frame instance segmentation, of which only anonymous boolean masks are exposed, or from SAM~2.1 (Hiera Small, automatic mask generation)~\cite{sam2_2025,sam21release}; each covers at least 196 pixels, at most 64 per frame.
A fragment carries the mean DINOv2 ViT-B/14 patch token in its mask~\cite{dinov2_2024}, passed through a frozen 128-dimensional projection trained once per mask source on 30 development houses disjoint from all training, selection and test houses, and its back-projected points, centroid and box.
Whether an entity should be visible, and how much of its last seen surface is now seen through, is computed per arm from the public depth and $M_{t-1}$ by a shared function.
Instance masks remove segmentation error unrelated to memory but expose no identity; SAM~2.1 masks test whether conclusions hold with a learned segmenter.

\subsection{Recall and frame-level joint assignment}
For each fragment $f$, recall unites a local channel (top $k{=}5$ entities by cosine within 3\,m) and a global channel (top $k'{=}3$, any distance, any state); eligibility does not depend on the entity state, so \arm{AssocOnly} receives the same far-away candidates as VSMT-lean.
The cost matrix has a row per fragment and a column per recalled entity plus one \op{BIRTH} column per fragment,
\begin{equation}
C[f,e]=-a(f,e),\qquad C[f,\mathrm{birth}_f]=-b(f),
\label{eq:cost}
\end{equation}
with association and birth logits $a$, $b$ and a forbidding cost for unrecalled pairs.
A minimum-cost rectangular assignment~\cite{kuhn1955}, canonicalised to the lexicographically smallest optimum, is compiled into \op{BIND} (active entity), \op{REACTIVATE} (dormant or retracted entity) or \op{BIRTH}: it decides in one step who is whom, who is new and who is back, and the learned heads only supply the costs.
For every unassigned, should-be-visible active or dormant entity, the existence head then decides \op{RETRACT} or \op{NOOP}.

\subsection{Learned cost heads}
The association head $a(f,e)$ reads 14 features (descriptor cosines, centroid distance, box IoU, size ratio, time since last seen, missed opportunities, state, cosine rank and margin, mutual best match, support-height difference); the existence head $r(e)$ reads 17 (visible fraction, free-space coverage of the last seen surface, viewing distance and angle, counts, time since last seen, best cosine to a current fragment and whether it is unassigned, state, and five history counters of see-through runs at coverage $\geq 0.70$ and $\geq 0.85$, matches since birth, judgeable frames and accumulated coverage); the birth head $b(f)$ reads 4.
Each head standardises its fields with training-house statistics (counts log-transformed) and applies a two-layer, 128-wide GELU MLP: 54,787 parameters in all (\arm{AssocOnly}: 35,842).
The loss is a softmax cross-entropy over $[\text{recalled columns},\mathrm{birth}_f]$ per labelled fragment plus a binary cross-entropy for existence with absent rows weighted by $w$, the present-to-absent ratio in the training records; the corrected score $\sigma(r(e)-\ln w)$ estimates how likely the entity is no longer where the memory keeps it, adjusted for how rare that case is in training, and the entity is retracted when the score reaches $\tau_r$.
Training uses AdamW (weight decay $10^{-4}$, one frame per batch, 20 epochs, learning rate cosine-decayed from $10^{-3}$ to $10^{-5}$, gradient clipping 1.0) and five seeds, in two rounds because the memory a head sees depends on its own decisions~\cite{dagger2011}: round~0 on records of the persistence-filter arm \arm{ELU-P}, round~1 adding records of the round-0 heads.
Checkpoints are scored on held-out training houses, never on validation: round~0 by total loss; in round~1, used for all results, the association and birth heads by association loss and the existence head by existence loss.

\subsection{Hindsight teacher and the candidate-before-teacher seal}
\label{sec:teacher}
Labels come from private instance truth at time $t$, not from future observations.
A fragment's dominant object is read from the instance map, and an entity's identity is the strict majority of the objects in its evidence (otherwise ambiguous).
The association target is the entity carrying the fragment's object, or \op{BIRTH} if none exists; the existence label is ``absent'' when the object has left the scene or the entity is no longer at its object by the node rule (Section~\ref{sec:metrics}).
Recall sets and association and birth features are sealed (stage A) before any private file is opened, and assignments and existence features after the solve (stage B); private files open only against a receipt naming both digests, and the teacher never creates, removes or reorders candidates.
Every decision falls into exactly one class: \emph{recall miss} (correct entity not recalled), \emph{teacher error} (ambiguous label), \emph{amortization error} (clear label, candidate present, different choice), correct, unlabelled, or duplicate (a second fragment of an object already labelled in the frame, which no student could satisfy since an entity takes one fragment per frame).
The ledger shows where errors arise; it does not by itself make a method better.

\subsection{Ablations and rule-based arms}
\label{sec:arms}
All arms share the front-end cache, recall, deduplication, executor, evaluator and, where they have a dormant state, the dormancy rule; only costs and existence decisions differ.
\arm{AssocOnly} retrains the same association head and solver with the same recipe on its own trajectories, with \op{BIND} and \op{BIRTH} only; it is the counterfactual for the lifecycle vocabulary.
\arm{NoVersion} uses VSMT-lean's heads but deletes retracted entities, which can then no longer be recalled or reactivated.
\arm{HeuristicLabel} trains the same network with the same recipe on \arm{ELU-P}'s public decisions instead of instance truth (\arm{ELU-P} fits three rates on training truth).
\arm{HandCost} keeps VSMT-lean's solver, vocabulary and executor with hand-written costs: descriptor cosine, a constant birth cost, and the current frame's free-space coverage for existence.
\arm{TAF} associates and fuses by thresholds as in ConceptGraphs~\cite{conceptgraphs2024} and never retracts.
\arm{ELU-P} adds to \arm{TAF} an existence log-odds that falls with reliable free-space coverage and rises with matches, in the style of Fusion++ and Perpetua~\cite{fusionpp2018,perpetua2025}; it retracts past a threshold and reactivates on a match.
\arm{RAC} adds to \arm{TAF} render-and-compare removal in the style of DSG~\cite{dsg2026}: repeated see-through evidence on the entity's projected surface retracts it.
\arm{LOW} lets the last observation win: the nearest entity within a distance takes it, otherwise birth; it never retracts.
These adapters reimplement published mechanisms on our interface and are not the original systems.
